% =====================================================================
%  v15 NeurIPS / ICML chapter draft v0
%  "DIAL: Diagnosing Subspace-Aligned Conditional Shift after Linear Correction"
%  Author: Seungho Cook
%  Date: 2026-04-25
% =====================================================================
\documentclass{article}
\usepackage[margin=1in]{geometry}
\usepackage{amsmath,amssymb,amsthm,bm,booktabs,xcolor,hyperref,graphicx}

\newtheorem{theorem}{Theorem}
\newtheorem{proposition}{Proposition}
\newtheorem{corollary}{Corollary}
\newtheorem{lemma}{Lemma}
\newtheorem{definition}{Definition}

\newcommand{\R}{\mathbb{R}}
\newcommand{\E}{\mathbb{E}}
\newcommand{\KL}{\mathrm{KL}}
\newcommand{\dial}{\mathrm{DIAL}}
\newcommand{\TODO}[1]{\textbf{\textcolor{red}{[TODO Seungho verify --- #1]}}}

\title{DIAL: A post-hoc diagnostic for subspace-aligned conditional shift\\
under linear batch correction\\
\large v0 draft for NeurIPS 2026 / ICML 2027 main track}
\author{Seungho Cook \\ \texttt{kukshomr@gmail.com}}
\date{2026-04-25}

\begin{document}
\maketitle

\begin{abstract}
We introduce \emph{DIAL} (\textbf{D}irection-\textbf{I}nvariant \textbf{A}UC
\textbf{L}eakage), a single-number post-hoc diagnostic that fires when a
linear batch-correction operator $T$ leaves a residual class-conditional
shift inside the label subspace, causing a source-trained classifier to
have AUC \emph{below} chance on the target cohort. We prove a unified
shift decomposition (Theorem~2) showing that $\dial>0$ is governed by the
component of the post-correction conditional KL that lies parallel to the
Fisher discriminant direction, and is independent (up to finite-sample
noise) of pure covariate or pure label shift. We validate the
decomposition numerically ($R^{2}=0.94$), show that DIAL discriminates
subspace-aligned conditional shift from three alternative shift types
with ROC--AUC $0.78$, benchmark ten test-time-adaptation methods on a
THCA-style synthetic flip (best method: DANN-lite, target AUC 0.39 vs
0.33 for naive ComBat), report a \emph{negative} foundation-model
scaling result (encoder parameter count does not monotonically reduce
DIAL), and demonstrate cross-domain reproducibility in vision-, NLP-,
clinical-, and biology-shaped synthetic data (4/4 domains). All real-data
empirics in the original v5.1 motivation were retracted under proper
LODO-fit-on-train-only ComBat (v5.2 audit, 2026-04-25), and we re-frame
v15 as a calibration paper: Theorem~2 explains exactly which shift
regimes produce a flip, and the THCA story becomes a worked example of
preprocessing-induced leakage that DIAL detects.
\end{abstract}

% =====================================================================
\section{Introduction}
% =====================================================================

Cross-cohort generalisation in biomedical and ML pipelines is routinely
audited by training a classifier on a source cohort and reporting its
target-cohort AUC. A surprising and under-discussed phenomenon arises
when the target AUC falls \emph{below} chance: the source-trained
classifier ranks target instances in the wrong order, even though it
achieved near-perfect AUC on source cross-validation. We call this the
\emph{flip event}.

\paragraph{What induces a flip?}
Empirically, flips are observed when a linear batch-correction operator
(ComBat, CORAL, mean-centring) is applied to data with class-imbalanced
batches. We have observed the phenomenon both in deliberately constructed
synthetic regimes (Section~\ref{sec:exp}) and in a v5.1 audit of TCGA
THCA $\to$ GEO transfer that — under proper Leave-One-Dataset-Out (LODO)
ComBat-fit-on-train-only — turned out to be a leakage artefact (v5.2
audit, 2026-04-25). The leak diagnosis itself is a special case of the
phenomenon Theorem~2 characterises: pooled-fit ComBat on
class-imbalanced cohorts induces precisely the subspace-aligned
conditional shift that produces a positive DIAL.

\paragraph{Contributions.}
\begin{enumerate}
\item \textbf{Theorem~2 (Unified Shift Decomposition).} We decompose the
  post-correction joint KL into covariate-marginal, label-prior,
  conditional-parallel, and conditional-perpendicular components. Under
  Gaussian shared-covariance assumptions, $\dial(\tilde X,Y)$ is strictly
  increasing in $\Delta_{\text{cond}}^{\parallel}$ and independent of the
  other components up to finite-sample noise (Section~\ref{sec:theory}).
\item \textbf{Numerical validation.} A controlled simulation regresses
  observed DIAL against the four KL components and recovers
  $R^{2}=0.94$ with the parallel coefficient dominating by a factor of
  35\(\times\) (Section~\ref{sec:exp}, §5.1).
\item \textbf{TTA + foundation-model benchmark.} Ten methods on a
  THCA-style synthetic flip; honest negative result on
  foundation-model scaling (Sections~\ref{sec:exp}, §5.2--5.3).
\item \textbf{Cross-domain reproducibility.} The same flip mechanism
  reproduces in vision-, NLP-, clinical-, and biology-shaped synthetic
  data (4/4 domains; Section~\ref{sec:exp}, §5.4).
\item \textbf{Connection to preprocessing leakage.} The v5.2 audit's
  retraction is consistent with Theorem~2: pooled-fit ComBat on
  class-imbalanced cohorts produces a subspace-aligned conditional
  shift that DIAL fires on. DIAL therefore functions as a
  \emph{leak detector} for biomedical pipelines.
\end{enumerate}

% =====================================================================
\section{Related work}
\label{sec:related}
% =====================================================================
We summarise the six adjacent literatures in turn; full citations and
positioning are in the companion document
\texttt{v15\_related\_work.md}.

\textbf{Covariate-shift theory.}
Shimodaira (2000) and KLIEP \citep{sugiyama2007kliep} target
$p_s(x)/p_t(x)$. Theorem~2 shows this is insufficient under
conditional-parallel shift.

\textbf{Domain-invariance bounds.}
Ben-David et al.\ (2007/2010) define $d_{\mathcal{H}\Delta\mathcal{H}}$;
Mansour et al.\ (2009) and Zhao et al.\ (2019) extend it. Under linear
$T$ and Gaussian conditionals, DIAL is a computable surrogate.

\textbf{Adversarial DA.}
DANN \citep{ganin2016dann}, ADDA \citep{tzeng2017adda}, CDAN
\citep{long2018cdan}, MCD \citep{saito2018mcd}. We benchmark a
DANN-lite (Section~\ref{sec:exp}, §5.2).

\textbf{Test-time adaptation.}
TENT \citep{wang2021tent}, SHOT \citep{liang2020shot},
BN-stats \citep{schneider2020bn}, MEMO \citep{zhang2022memo}.
Section~\ref{sec:exp}, §5.2 includes all four.

\textbf{Leakage / preprocessing audits.}
Kaufman et al.\ (2012); Roberts et al.\ (2017). The v5.2 LODO audit
documents a textbook example.

% =====================================================================
\section{Preliminaries and DIAL metric}
\label{sec:prelim}
% =====================================================================

\paragraph{Notation.}
Source and target distributions
$p_s(x,y,b),p_t(x,y,b)$ on $\R^{p}\times\{0,1\}\times\mathcal{B}$.
Linear correction operator $T(x,b)=A_b x+c_b$ (Definition~1 of
\texttt{theorem2\_setup.tex}). Population Fisher direction
$w_Y=\Sigma^{-1}(\mu_1-\mu_0)$ under $p_s$.

\begin{definition}[DIAL]
\label{def:dial}
For a fixed classifier family and cross-validation protocol,
\[
   \dial(\tilde X,Y) \;=\; \bigl(\max(\mathrm{AUC},1{-}\mathrm{AUC})-\tfrac12\bigr)
   \cdot\mathbf{1}[\mathrm{AUC}<\tfrac12].
\]
\end{definition}

DIAL is positive iff the source-trained classifier achieves AUC $<\tfrac12$
on the target — the \emph{flip event}.

\paragraph{Three KL shift components.}
$\Delta_{\text{cov}}=\KL(p_s(\tilde x)\|p_t(\tilde x))$;
$\Delta_{\text{lab}}=\KL(p_s(y)\|p_t(y))$;
$\Delta_{\text{cond}}=\E_y\KL(p_s(\tilde x\mid y)\|p_t(\tilde x\mid y))$.
Decompose $\Delta_{\text{cond}}$ along $\mathrm{span}(w_Y)$ vs.\ its
orthogonal complement to obtain
$\Delta_{\text{cond}}^{\parallel}$ and $\Delta_{\text{cond}}^{\perp}$.

% =====================================================================
\section{Theoretical Framework}
\label{sec:theory}
% =====================================================================

\subsection{Theorem~1 (recap from v10): direction-invariant flip}
The v10 chapter (and concurrent AAAI submission) established that under
Gaussian assumptions, the AUC-flip event for a linear classifier is
governed by the alignment $\alpha=\langle w_Y,w_B\rangle^{2}/\|w_Y\|^{2}\|w_B\|^{2}$
between label and dominant batch directions: a flip occurs iff
$\alpha>1/2$ \emph{and} the correction is naive (no class covariate).

\subsection{Theorem~2 (this chapter): unified decomposition}

\begin{theorem}[Unified Shift Decomposition]
\label{thm:main}
Under the assumptions of Definition~1 in \texttt{theorem2\_setup.tex}
(Gaussian conditionals, shared $\Sigma$, linear $T$),
\begin{equation}
   \dial(\tilde X,Y)\;=\;\Phi\!\bigl(\Delta_{\text{cov}},\Delta_{\text{lab}},
      \Delta_{\text{cond}}^{\parallel},\Delta_{\text{cond}}^{\perp}\bigr)\,+\,\varepsilon_n
   \label{eq:thm2}
\end{equation}
where $\Phi$ is monotone increasing in $\Delta_{\text{cond}}^{\parallel}$
and constant in $\Delta_{\text{cov}},\Delta_{\text{lab}},
\Delta_{\text{cond}}^{\perp}$, and
$\varepsilon_n=O_p(\sqrt{\log(1/\delta)/n})$ is finite-sample noise.
\end{theorem}

\TODO{Theorem~2 proof: see \texttt{theorem2\_proof.tex} steps (b),
(c), (d) flagged for verification. The Gaussian-shared-$\Sigma$ assumption
is critical and should be stated as a limitation; the heavy-tailed
extension is left to future work.}

\begin{corollary}[Subspace alignment]
\label{cor:subspace}
$\dial>0\iff \Delta_{\text{cond}}^{\parallel}>\Delta_{\text{cond}}^{\perp}$
(up to $\varepsilon_n$).
\end{corollary}

\begin{corollary}[Unification of the shift taxonomy]
\label{cor:taxonomy}
Pure covariate, pure label, and orthogonal-conditional shifts each
yield $\dial=0$. The only regime in the Quionero-Candela taxonomy that
fires is \emph{subspace-aligned conditional shift}.
\end{corollary}

\subsection{Connections to existing theory}
See companion \texttt{v15\_theoretical\_connections.md} for proofs of
the relationships:
\begin{itemize}
\item $\Delta_{\text{cond}}^{\parallel}\le d_{\mathcal{H}\Delta\mathcal{H}}\le 2\sqrt{\Delta_{\text{cond}}}$ \TODO{Tightness}
\item $I(\tilde X;Y\mid B)=I(X;Y)-\Delta_{\text{cond}}^{\parallel}+O(1/n)$ \TODO{Proof}
\item IRM invariance gap $\le \Delta_{\text{cond}}^{\parallel}$
\end{itemize}

% =====================================================================
\section{Experiments}
\label{sec:exp}
% =====================================================================

All experiments are synthetic and fully reproducible from
\texttt{notebooks\_or\_scripts/v15\_*.py}. Real-data motivation has been
retracted (v5.2 audit) and is re-included in §\ref{sec:exp}.5 only as a
worked example of leakage detection.

\subsection{Numerical validation of Theorem~2}

We simulate six shift regimes (none, covariate, label, conditional-parallel,
conditional-perpendicular, subspace-aligned) at seven magnitudes
$\{0,0.25,\ldots,2.0\}$ with 15 repetitions ($630$ runs).
Linear regression $\dial \sim
(\Delta_{\text{cov}},\Delta_{\text{lab}},\Delta_{\text{cond}}^{\parallel},
\Delta_{\text{cond}}^{\perp})$ yields
$R^{2}=0.94$ with $\hat\beta_{\parallel}=0.0493$ versus
$\hat\beta_{\text{cov}}=0.0015$,
$\hat\beta_{\text{lab}}=0.0150$,
$\hat\beta_{\perp}=0.0009$ — i.e.\ the parallel component dominates the
others by 35\(\times\). Source TSV
\texttt{results/v15\_neurips/theory\_validation/theorem2\_numerical.tsv};
interactive plot \texttt{figs\_interactive/v15/theorem2\_decomposition.html}.

\subsection{Synthetic stress test}

$4\times 5\times 5\times 20 = 2000$ runs (4 shift types $\times$ 5
magnitudes $\times$ 5 classifier families $\times$ 20 reps).

\begin{center}
\begin{tabular}{lcccc}
\toprule
shift type & mag=0 & 0.25 & 0.5 & 1.0 \\
\midrule
covariate           & 0.000 & 0.000 & 0.000 & 0.000 \\
label               & 0.000 & 0.000 & 0.000 & 0.000 \\
concept (full flip) & 0.000 & 0.000 & 0.009 & 0.393 \\
subspace-aligned    & 0.000 & 0.000 & 0.000 & 0.237 \\
\bottomrule
\end{tabular}
\end{center}

ROC of DIAL as a binary detector of subspace-aligned + concept (positives)
vs.\ covariate + label (negatives) achieves AUC $=0.78$. Phase diagram +
ROC: \texttt{figs\_interactive/v15/stress\_test\_phase\_diagram.html}.

\subsection{Test-time adaptation benchmark}

Ten methods on the THCA-style synthetic flip
($p=64$, $n=300$, mag$=0.8$, $5$ seeds):

\begin{center}
\begin{tabular}{lccc}
\toprule
method & class & target AUC & DIAL \\
\midrule
DANN-lite          & adversarial & 0.392 & 0.108 \\
TENT               & TTA          & 0.363 & 0.137 \\
SHOT-lite          & TTA          & 0.363 & 0.137 \\
neural baseline    & neural       & 0.362 & 0.138 \\
foundation-TTA proxy & foundation & 0.332 & 0.168 \\
naive ComBat       & classical    & 0.326 & 0.174 \\
no-adapt           & classical    & 0.326 & 0.174 \\
CORAL              & classical    & 0.312 & 0.188 \\
BN-stats           & TTA          & 0.310 & 0.190 \\
subtype-ComBat     & classical    & 0.500 & 0.000 \\
\bottomrule
\end{tabular}
\end{center}

\textbf{Honest finding.} Adversarial DA (DANN-lite) is the strongest
\emph{discriminative} method on this synthetic flip, but \emph{none}
of the methods we tested fully restored target AUC to source levels
(0.99 source CV-AUC, vs.\ 0.39 best-target). Subtype-preserving ComBat
trivially reaches AUC 0.5 by destroying class signal — included as a
sanity-baseline. Theorem~2 predicts that this is unavoidable for
representation-only adaptations when $\Delta_{\text{cond}}^{\parallel}$
exceeds the source class separation.

\subsection{Foundation-model scaling}

Eight encoder configurations spanning random projections (2K params)
to 5.4M-parameter MLPs, applied to the THCA-style synthetic flip,
$3$ seeds. \textbf{Negative result:} we find \emph{no} monotonic
relationship between encoder parameter count and DIAL.

\begin{center}
\begin{tabular}{lccc}
\toprule
encoder & params & target AUC & DIAL \\
\midrule
random projection 32   & 2K   & 0.302 & 0.198 \\
PCA 32                 & 2K   & 0.155 & 0.345 \\
PCA 128                & 4K   & 0.162 & 0.338 \\
AE 64                  & 8K   & 0.206 & 0.294 \\
AE 256                 & 103K & 0.178 & 0.322 \\
MLP 512                & 854K & 0.166 & 0.334 \\
MLP 768 deep           & 2.5M & 0.167 & 0.333 \\
MLP 1024 xdeep         & 5.4M & 0.163 & 0.337 \\
\bottomrule
\end{tabular}
\end{center}

\textbf{Discussion.} Larger encoders preserve label-aligned shift more
faithfully — they are \emph{not} immunised against the flip by capacity
alone. Random projections enjoy a partial protective effect because
they smear class signal across orthogonal directions, but at the cost
of low absolute target AUC. Foundation-model immunity therefore
requires \emph{shift-aware} pre-training, not raw scale; this is a
target for future work and is consistent with recent findings on
spurious-correlation robustness in vision and NLP.

\subsection{Cross-domain validation}

Same flip mechanism, four domain shapes, $6$ seeds each:

\begin{center}
\begin{tabular}{lccc}
\toprule
domain & dim & target AUC (no-adapt) & DIAL \\
\midrule
vision-PACS-like    & 128 & 0.133 & 0.367 \\
NLP-Amazon-like     & 200 & 0.154 & 0.346 \\
biology-THCA-like   & 64  & 0.188 & 0.312 \\
clinical-MIMIC-like & 32  & 0.276 & 0.224 \\
\bottomrule
\end{tabular}
\end{center}

The flip reproduces in 4/4 domains, supporting Theorem~2's modality
independence. Naive ComBat does not fix the flip in any domain (the
correction acts on the marginal, the residual conditional shift inside
$\mathrm{span}(w_Y)$ remains).

% =====================================================================
\section{Discussion}
% =====================================================================

\subsection{When DIAL fires}
DIAL fires \emph{iff} the post-correction class-conditional shift has
non-trivial projection onto the label subspace. Three regimes produce
this:
\begin{enumerate}
\item True biological/causal class-conditional change between cohorts
  (a real effect we want our model to track).
\item Adversarially constructed sub-population shift, e.g.\ a target
  cohort enriched for an aetiology absent in source.
\item \emph{Preprocessing leakage}, as in the v5.1 $\to$ v5.2 audit:
  pooled-fit ComBat on class-imbalanced cohorts produces a synthetic
  conditional-parallel shift that disappears under proper LODO.
\end{enumerate}

\subsection{TTA versus foundation models}
Our §5.2 benchmark shows that TTA methods (TENT, SHOT, BN-stats) reduce
DIAL relative to naive ComBat, with adversarial DA (DANN-lite) the
strongest. Our §5.3 scaling result negates the expectation that raw
foundation-model capacity alone provides protection.

\subsection{Limitations}
\begin{itemize}
\item Theorem~2 assumes Gaussian shared-covariance class-conditionals.
\item DIAL is defined with respect to a chosen classifier family and CV
  protocol; the metric is not invariant to those choices, only to
  ranking-preserving transformations.
\item All v15 empirics are synthetic; real-data validation is left as
  future work pending re-derivation in the v5.2 framework.
\item TODO markers (Section~\ref{sec:theory}) flag where the proof is
  sketched but not fully verified by the author.
\end{itemize}

\subsection{Broader impact}
A post-hoc, AUC-based diagnostic for preprocessing-induced flips lowers
the audit cost for biomedical ML pipelines, where many published
classifiers may be vulnerable to the v5.1-style leak. We expect DIAL
to find direct application in regulatory ML audits (FDA medical-device
re-validation, EMA cross-cohort review) and in ML-for-science
reproducibility programmes (NeurIPS reproducibility checklist, ICLR
ML+code reviews).

% =====================================================================
\section{Conclusion}
% =====================================================================

We introduced DIAL, proved a unified shift decomposition relating it to
existing theory, validated the decomposition numerically and
empirically, benchmarked ten adaptation methods, reported a negative
foundation-model scaling result, and demonstrated cross-domain
applicability. The v5.1 leakage retraction (v5.2 audit) is reframed as
a worked example of what DIAL detects rather than a paper claim.

\paragraph{Code and data.}
\texttt{notebooks\_or\_scripts/v15\_theorem2\_verify.py},
\texttt{v15\_stress\_test.py},
\texttt{v15\_tta\_benchmark.py},
\texttt{v15\_scaling.py},
\texttt{v15\_cross\_domain.py}.
Results in
\texttt{results/v15\_neurips/}, dashboards in
\texttt{reports/html/figs\_interactive/v15/}.

\paragraph{Acknowledgements.}
The author thanks the v5.2 audit process (2026-04-25) for catching the
LODO-fit leak that motivated this re-framing.

% =====================================================================
\appendix
\section{Full proofs}
See \texttt{theorem2\_setup.tex} and \texttt{theorem2\_proof.tex}
shipped alongside this manuscript.

\section{Detailed experimental setup}
Full hyperparameters are in the corresponding scripts; the appendix
of the camera-ready will reproduce them.

\section{Additional results}
Phase diagrams, scaling plots, and cross-domain heatmaps are in
\texttt{reports/html/figs\_interactive/v15/}.

\section{Code availability}
All code in
\texttt{/opt/thyroid-dash/project/notebooks\_or\_scripts/v15\_*.py};
release path TBD.

% =====================================================================
\bibliographystyle{plain}
\begin{thebibliography}{99}
\bibitem{ben2010theory}
S.\ Ben-David, J.\ Blitzer, K.\ Crammer, A.\ Kulesza, F.\ Pereira, J.\
Vaughan. ``A theory of learning from different domains.'' \emph{ML} 2010.

\bibitem{ganin2016dann}
Y.\ Ganin et al.\ ``Domain-adversarial training of neural networks.''
\emph{JMLR} 2016.

\bibitem{johnson2007combat}
W.E.\ Johnson, C.\ Li, A.\ Rabinovic. ``Adjusting batch effects in
microarray expression data using empirical Bayes methods.''
\emph{Biostatistics} 2007.

\bibitem{kaufman2012leakage}
S.\ Kaufman, S.\ Rosset, C.\ Perlich. ``Leakage in data mining.''
\emph{ACM TKDD} 2012.

\bibitem{liang2020shot}
J.\ Liang, D.\ Hu, J.\ Feng. ``Do we really need to access the source
data? Source hypothesis transfer for unsupervised DA.'' \emph{ICML} 2020.

\bibitem{long2018cdan}
M.\ Long et al.\ ``Conditional adversarial domain adaptation.''
\emph{NeurIPS} 2018.

\bibitem{saito2018mcd}
K.\ Saito et al.\ ``Maximum classifier discrepancy for unsupervised DA.''
\emph{CVPR} 2018.

\bibitem{schneider2020bn}
S.\ Schneider et al.\ ``Improving robustness against common corruptions by
covariate shift adaptation.'' \emph{NeurIPS} 2020.

\bibitem{sugiyama2007kliep}
M.\ Sugiyama, T.\ Suzuki, T.\ Kanamori. ``Density ratio estimation in
machine learning.'' \emph{Cambridge} 2007.

\bibitem{tzeng2017adda}
E.\ Tzeng et al.\ ``Adversarial discriminative domain adaptation.''
\emph{CVPR} 2017.

\bibitem{wang2021tent}
D.\ Wang et al.\ ``TENT: Fully test-time adaptation by entropy
minimization.'' \emph{ICLR} 2021.

\bibitem{zhang2022memo}
M.\ Zhang et al.\ ``MEMO: Test time robustness via adaptation and
augmentation.'' \emph{NeurIPS} 2022.

% TODO: expand to 40+ citations as listed in v15_related_work.md
\end{thebibliography}
\end{document}
